% ============================================================================
% Chapter 32: Charts and Data Management
% ============================================================================
\chapter{Charts and Data Management}
\label{ch:excel-charts}

\begin{objectives}
  \item Create bar charts, line charts, and pie charts from data
  \item Sort data in ascending and descending order
  \item Use basic filters to find specific data
  \item Choose the right chart type for your data
\end{objectives}

\bytetip[tip]{Charts turn boring numbers into beautiful visuals. A single chart can tell a story that would take paragraphs of text. Let's learn to create stunning charts!}

\section{Types of Charts}
\label{sec:chart-types}
\index{MS Excel!charts}

% --- Figure 32.1: Chart Type Guide ---
\begin{figure}[H]
\centering
\begin{tikzpicture}
  % Bar chart panel
  \node[draw=ModuleBlue, fill=ModuleBlue!5, rounded corners=6pt, 
        minimum width=4cm, minimum height=4cm, align=center] (bar) at (0,0) {};
  \begin{scope}[shift={(-1.2,-1)}]
    \fill[ModuleBlue!60] (0,0) rectangle (0.5,1.5);
    \fill[FunCoral!60] (0.7,0) rectangle (1.2,2.0);
    \fill[LabGreen!60] (1.4,0) rectangle (1.9,1.0);
  \end{scope}
  \node[font=\small\bfseries\sffamily, text=ModuleBlue, above=0.1cm] at (0,2) {Bar Chart};
  \node[font=\tiny\sffamily, text=NavyBlue, below=0.1cm] at (0,-2) {Compare categories};
  
  % Line chart panel
  \node[draw=FunPurple, fill=FunPurple!5, rounded corners=6pt,
        minimum width=4cm, minimum height=4cm] (line) at (5.5,0) {};
  \begin{scope}[shift={(4,-0.8)}]
    \draw[FunPurple, line width=1.5pt] (0.3,0) -- (0.8,0.8) -- (1.3,0.5) -- (1.8,1.5) -- (2.3,1.2);
    \foreach \x/\y in {0.3/0, 0.8/0.8, 1.3/0.5, 1.8/1.5, 2.3/1.2}{
      \fill[FunPurple] (\x,\y) circle (3pt);
    }
  \end{scope}
  \node[font=\small\bfseries\sffamily, text=FunPurple, above=0.1cm] at (5.5,2) {Line Chart};
  \node[font=\tiny\sffamily, text=NavyBlue, below=0.1cm] at (5.5,-2) {Show trends over time};
  
  % Pie chart panel
  \node[draw=LabGreen, fill=LabGreen!5, rounded corners=6pt,
        minimum width=4cm, minimum height=4cm] (pie) at (11,0) {};
  \begin{scope}[shift={(11,0)}]
    \fill[FunSky!70] (0,0) -- (0,1) arc (90:210:1) -- cycle;
    \fill[FunCoral!70] (0,0) -- (-0.87,-0.5) arc (210:330:1) -- cycle;
    \fill[FunYellow!80] (0,0) -- (0.87,-0.5) arc (330:450:1) -- cycle;
  \end{scope}
  \node[font=\small\bfseries\sffamily, text=LabGreen, above=0.1cm] at (11,2) {Pie Chart};
  \node[font=\tiny\sffamily, text=NavyBlue, below=0.1cm] at (11,-2) {Parts of a whole};
\end{tikzpicture}
\caption{Three chart types and when to use each one}
\label{fig:chart-types}
\end{figure}

\begin{labstep}[Creating a Chart]
\begin{enumerate}[label=\protect\stepcircle{\arabic*}, leftmargin=2cm, itemsep=3pt]
  \item Select the data you want to chart (including headers).
  \item Go to \textbf{Insert} tab $\rightarrow$ \textbf{Charts} group.
  \item Choose your chart type (Bar, Line, or Pie).
  \item Click the chart to customise: add a title, labels, and colours.
  \item Resize and move the chart as needed.
\end{enumerate}
\end{labstep}

\section{Sorting Data}
\label{sec:sorting}
\index{MS Excel!sorting}

% --- Figure 32.2: Sorting Illustrated ---
\begin{figure}[H]
\centering
\begin{tikzpicture}
  % Before sorting
  \node[font=\small\bfseries\sffamily, text=NavyBlue] at (2,4) {Before};
  \foreach \y/\name in {3/Alice, 2/Zara, 1/Bob, 0/Mia}{
    \draw[NavyBlue!30, line width=0.5pt] (0.5,\y) rectangle (3.5,\y+0.8);
    \node[font=\small\sffamily] at (2,\y+0.4) {\name};
  }
  
  % Arrow
  \draw[LabGreen, line width=2pt, ->] (4.2,2) -- (5.8,2);
  \node[above, font=\scriptsize\sffamily\bfseries, text=LabGreen] at (5,2.1) {Sort A$\rightarrow$Z};
  
  % After sorting
  \node[font=\small\bfseries\sffamily, text=NavyBlue] at (8,4) {After};
  \foreach \y/\name in {3/Alice, 2/Bob, 1/Mia, 0/Zara}{
    \draw[NavyBlue!30, line width=0.5pt] (6.5,\y) rectangle (9.5,\y+0.8);
    \fill[FunMint!15] (6.5,\y) rectangle (9.5,\y+0.8);
    \node[font=\small\sffamily] at (8,\y+0.4) {\name};
  }
\end{tikzpicture}
\caption{Sorting rearranges data in alphabetical or numerical order}
\label{fig:sorting}
\end{figure}

To sort: Select your data $\rightarrow$ \textbf{Data} tab $\rightarrow$ click \faSortAlphaDown\ \textbf{Sort A to Z} or \faSortAlphaDown\ \textbf{Sort Z to A}.

\bytetip[fun]{Excel can create over \textbf{20 different chart types}, including radar charts, waterfall charts, treemaps, and even 3D surface plots! For school, you'll mostly use bar, line, and pie charts.}

\begin{funfact}
The first pie chart was created in 1801 by \textbf{William Playfair}, a Scottish engineer. He used it to show the proportions of the Turkish Empire's territory in Asia, Europe, and Africa. Pie charts are over 200 years old!
\end{funfact}

\begin{reviewqs}
  \item Name three types of charts you can create in Excel.
  \item When would you use a pie chart instead of a bar chart?
  \item What are the steps to create a chart from your data?
  \item How do you sort data alphabetically?
  \item What is a filter and how does it help?
\end{reviewqs}

\begin{challenge}
Enter marks for 5 students in 4 subjects. Create a \textbf{bar chart} comparing all students' total marks. Then create a \textbf{pie chart} showing how one student's marks are distributed across subjects. Add chart titles and labels.
\end{challenge}
